\documentclass[11pt]{article}
\usepackage[preprint]{acl}
\usepackage[T1]{fontenc}
\usepackage[utf8]{inputenc}
\usepackage{times}
\usepackage{latexsym}
\usepackage{microtype}
\usepackage{amsmath}
\usepackage{booktabs}
\usepackage{tabularx}
\usepackage{multirow}
\usepackage{graphicx}
\usepackage{enumitem}
\usepackage{hyperref}
\renewcommand{\ttdefault}{cmtt}

\title{Multimodal Hateful Meme Detection: Mid-Project Report}

\author{Lexical Leaders \\
Mudit Gupta \and Viral Shah \and Adarsh Sharma \and Bhautik Ramani \\
Introduction to Natural Language Processing, March 2026 \\
\texttt{https://github.com/MuditGupta2502/Meta-Hateful-Meme-Detection}}

\begin{document}
\maketitle

\begin{abstract}
This mid-submission report presents our project direction and current progress on multimodal hateful meme detection. Our core idea is to detect harmful memes by jointly modeling image and text context with a frozen CLIP backbone and a trainable cross-attention fusion head. So far, we have completed the Phase 1 baseline and a stronger Phase 1 extension based on class-balanced sampling. On the balanced development split, the baseline reaches AUROC 0.7077 and the extension reaches AUROC 0.7306. We also completed a controlled ablation that validates cross-attention over late-fusion concatenation (+0.0399 AUROC). For the planned Phase 2, we will focus on deployment realism (OCR), selective encoder adaptation (LoRA/PEFT), and an ethical safety layer with human-in-the-loop intervention for uncertain or high-risk predictions.
\end{abstract}

\section{Project Idea and Motivation}
Hateful memes are difficult to detect because hateful intent usually emerges from \textit{interaction} between image and text, not from one modality in isolation. A seemingly harmless image can become hateful with a specific caption, and a provocative sentence can be benign depending on visual context. This makes the task fundamentally multimodal.

Our project aims to build a practical and reproducible multimodal detection system that is both computationally efficient and extensible. Instead of full end-to-end fine-tuning of very large encoders, we freeze a strong pretrained backbone and train a smaller fusion/classification component. This design is suitable for academic constraints (limited GPU budget) while still allowing methodical experimentation.

\textbf{Primary objective.}
Given meme image $v$ and text $t$, learn a classifier
\[
f(v,t) \rightarrow \hat{y} \in \{0,1\},
\]
where $1$ denotes hateful and $0$ denotes non-hateful.

\textbf{Mid-project objectives.}
\begin{itemize}[leftmargin=*,nosep]
  \item Build a stable multimodal baseline with frozen CLIP encoders and cross-attention fusion.
  \item Improve minority-class robustness under class imbalance.
  \item Validate architecture decisions using controlled ablations.
  \item Define a realistic second phase with safety and deployment-focused additions.
\end{itemize}

\section{Background and Related Work}
The Hateful Memes Challenge \cite{kiela2020hateful} formalized this task with ``benign confounders'' so that unimodal shortcuts fail by design. This motivates fusion architectures that explicitly reason over image-text interactions.

Pretrained vision-language models, especially CLIP \cite{radford2021clip}, provide strong aligned embeddings for image and text. However, direct zero-shot similarity scoring is typically insufficient for subtle hate detection because the decision boundary depends on dataset-specific context and label semantics.

Transformer attention \cite{vaswani2017attention} motivates our cross-attention fusion block, where each modality attends to the other before classification. For imbalanced classification, focal loss \cite{lin2017focal} provides hard-example emphasis that is particularly useful when minority errors are costly.

For the planned second phase, we consider parameter-efficient adaptation via LoRA \cite{hu2022lora} so that representation quality can improve without full encoder fine-tuning. We also include model-risk safeguards inspired by broader responsible-AI moderation practice \cite{weidinger2021ethical}.

\section{Dataset}
We use the Meta AI Hateful Memes dataset \cite{kiela2020hateful}. The dataset provides image paths, meme text, and binary labels for train/dev splits; the public test split is unlabeled.

\begin{table}[h]
\centering
\small
\caption{Dataset split statistics used in this project.}
\begin{tabular}{lrrrr}
\toprule
Split & Total & Hateful & Non-hateful & Ratio \\
\midrule
Train & 8,500 & 3,050 & 5,450 & 1.79:1 \\
Dev   & 500   & 250   & 250   & 1:1 \\
Test  & 1,000 & --    & --    & -- \\
\bottomrule
\end{tabular}
\label{tab:data}
\end{table}

\textbf{Observed properties from EDA.}
\begin{itemize}[leftmargin=*,nosep]
  \item Meme text is short (mean \(\approx 11.7\) words), with long-tail stylistic noise.
  \item The train split is moderately imbalanced, requiring explicit mitigation.
  \item The development split is balanced and therefore suitable for robust threshold-free ranking metrics (AUROC).
\end{itemize}

\subsection{Exploratory Data Analysis Summary}
We performed an initial EDA pass to verify whether training assumptions match the actual data distribution.

\begin{table}[h]
\centering
\small
\caption{Word-count statistics on the training split.}
\begin{tabular}{lrrrr}
\toprule
Class & Mean & Median & Max & 95th pct. \\
\midrule
Non-hateful & 11.2 & 10 & 69 & 23 \\
Hateful & 12.8 & 11 & 70 & 26 \\
Overall & 11.7 & 10 & 70 & 24 \\
\bottomrule
\end{tabular}
\label{tab:wordstats}
\end{table}

These numbers confirm that meme text length is comfortably inside CLIP token limits, so truncation is not a dominant error source.

\begin{table}[h]
\centering
\small
\caption{Top frequent content tokens (stop-words removed).}
\begin{tabular}{ll}
\toprule
Hateful cluster examples & Non-hateful cluster examples \\
\midrule
``white'', ``black'', ``muslim'' & ``love'', ``happy'', ``friend'' \\
``illegal'', ``parasite'', ``send'' & ``birthday'', ``dog'', ``cat'' \\
``terror'', ``country'', ``back'' & ``food'', ``work'', ``today'' \\
\bottomrule
\end{tabular}
\label{tab:topwordsmid}
\end{table}

The lexical separation is non-trivial but incomplete, reinforcing the need for multimodal context instead of text-only filtering.

\section{Proposed Methodology and Architecture}
\subsection{Model Overview}
Our baseline architecture has three stages:
\begin{enumerate}[leftmargin=*]
  \item \textbf{Frozen multimodal encoding:} CLIP ViT-B/32 encodes image and text into 512-d vectors.
  \item \textbf{Bidirectional cross-attention fusion:} image queries attend to text features and text queries attend to image features.
  \item \textbf{Classification head:} a lightweight MLP predicts hateful vs non-hateful.
\end{enumerate}

Only the fusion and classifier layers are trainable (about 2.4M parameters), while the CLIP backbone remains frozen (about 151M parameters).

\subsection{Training Design}
\textbf{Loss.} Weighted focal loss with class weights \([1.0,\,1.786]\) and \(\gamma=2\).

\textbf{Optimizer and schedule.} AdamW, learning rate \(10^{-4}\), cosine decay with linear warmup.

\textbf{Class imbalance strategy.}
\begin{itemize}[leftmargin=*,nosep]
  \item Baseline: weighted focal loss only.
  \item Phase 1 extension: weighted focal loss + \texttt{WeightedRandomSampler} to create near-balanced sampling per epoch.
\end{itemize}

\subsection{Ablation Methodology}
To test whether cross-attention itself is useful, we compare:
\begin{itemize}[leftmargin=*,nosep]
  \item \textbf{Late Fusion:} concatenation of frozen CLIP image/text embeddings + MLP.
  \item \textbf{Cross-Attention Fusion:} full proposed architecture.
\end{itemize}
All other training choices are held fixed.

\subsection{Implementation Footprint}
\begin{table}[h]
\centering
\small
\caption{Parameter and compute footprint (current implementation).}
\begin{tabular}{lrr}
\toprule
Component & Parameters & Trainable \\
\midrule
CLIP ViT-B/32 backbone & \(\sim\)151.3M & No \\
Cross-attention fusion block & \(\sim\)2.1M & Yes \\
Classification head (MLP) & \(\sim\)0.3M & Yes \\
\midrule
Total & \(\sim\)153.7M & \(\sim\)2.4M (1.6\%) \\
\bottomrule
\end{tabular}
\label{tab:params}
\end{table}

This parameter-efficiency is important for reproducibility on limited hardware (Kaggle T4 class GPUs) while still allowing architecture-level experimentation.

\subsection{Core Objective Functions}
The classifier is trained with weighted focal loss:
\[
\mathcal{L}_{\text{focal}} = -\alpha_t (1-p_t)^{\gamma}\log(p_t),
\]
where \(p_t\) is the model probability assigned to the true class, \(\gamma=2\), and \(\alpha_t\) follows class imbalance.

For a batch \(\mathcal{B}\), the optimization objective is:
\[
\min_\theta \; \frac{1}{|\mathcal{B}|} \sum_{(v_i,t_i,y_i)\in\mathcal{B}}
\mathcal{L}_{\text{focal}}\big(y_i, f_\theta(v_i,t_i)\big).
\]

In the cross-attention module, image token representations \(V\) attend to text tokens \(T\), and text attends back to image:
\[
\text{Attn}(Q,K,V)=\text{softmax}\left(\frac{QK^\top}{\sqrt{d_k}}\right)V.
\]
This bidirectional interaction helps the model capture meaning that only emerges from image-text pairing.

\section{Evaluation Options and Metrics}
We evaluate on the 500-sample balanced dev split.

\textbf{Primary metric:} AUROC (threshold-independent ranking quality).

\textbf{Secondary metrics:} Accuracy, F1 (hateful class), Macro-F1, confusion matrix.

\textbf{Why these metrics.}
\begin{itemize}[leftmargin=*,nosep]
  \item AUROC is robust to threshold choice and is widely used for this task setup.
  \item Macro-F1 captures class-balance quality and guards against one-class dominance.
  \item Confusion matrices support practical moderation analysis (false positives vs false negatives).
\end{itemize}

We report best checkpoints by AUROC and track metric trajectories across epochs for overfitting diagnosis.

\begin{table*}[t]
\centering
\small
\caption{Main hyperparameters for the completed mid-phase experiments.}
\begin{tabular}{ll}
\toprule
Category & Setting \\
\midrule
Backbone & \texttt{openai/clip-vit-base-patch32} (frozen) \\
Fusion module & Bidirectional cross-attention, 4 heads \\
Classifier head & 3-layer MLP with dropout \\
Loss & Weighted focal loss, \(\gamma=2\), class weights \([1.0,1.786]\) \\
Optimizer & AdamW \\
Learning rate & \(1\times 10^{-4}\) \\
Weight decay & \(1\times 10^{-4}\) \\
Scheduler & Cosine decay with linear warmup \\
Batch size & 32 \\
Phase 1 baseline & 10 epochs, warmup 100 \\
Phase 1 extension & 20 epochs, warmup 200, weighted random sampler \\
Primary checkpoint criterion & Validation AUROC \\
Runtime footprint & Approximately 35 minutes (10-epoch run, Kaggle T4) \\
\bottomrule
\end{tabular}
\label{tab:hyperparams}
\end{table*}

\section{Interim Experimental Results}
\subsection{Phase 1 Baseline (10 epochs)}
\begin{table*}[t]
\centering
\small
\caption{Phase 1 baseline: epoch-wise validation trajectory.}
\begin{tabular}{rccccc}
\toprule
Epoch & Train Loss & Val Loss & Val AUROC & Val F1 & Macro-F1 \\
\midrule
1 & 0.298 & 0.356 & 0.639 & 0.600 & 0.531 \\
2 & 0.266 & 0.346 & 0.677 & 0.644 & 0.611 \\
3 & 0.245 & 0.344 & 0.681 & 0.665 & 0.633 \\
4 & 0.229 & 0.354 & 0.695 & 0.650 & 0.629 \\
5 & 0.219 & 0.358 & 0.702 & 0.695 & 0.664 \\
6 & 0.207 & 0.359 & 0.706 & 0.689 & 0.670 \\
7 & 0.196 & 0.367 & 0.707 & 0.682 & 0.666 \\
8 & 0.186 & 0.381 & 0.703 & 0.673 & 0.664 \\
9 & 0.178 & 0.389 & 0.708 & 0.664 & 0.658 \\
10 & 0.170 & 0.395 & 0.705 & 0.655 & 0.653 \\
\bottomrule
\end{tabular}
\label{tab:phase1epochs}
\end{table*}

\subsection{Phase 1 Extension (20 epochs, balanced sampler)}
\begin{table*}[t]
\centering
\small
\caption{Phase 1 extension: epoch-wise validation trajectory (best AUROC at epoch 8).}
\begin{tabular}{rccccc}
\toprule
Epoch & Train Loss & Val Loss & Val AUROC & Val F1 & Macro-F1 \\
\midrule
1 & 0.301 & 0.314 & 0.6696 & 0.6918 & 0.5441 \\
2 & 0.251 & 0.318 & 0.6831 & 0.6875 & 0.5660 \\
3 & 0.229 & 0.319 & 0.7005 & 0.7018 & 0.6206 \\
4 & 0.208 & 0.342 & 0.7049 & 0.6985 & 0.6349 \\
5 & 0.188 & 0.345 & 0.7206 & 0.7018 & 0.6532 \\
6 & 0.164 & 0.408 & 0.7147 & 0.6991 & 0.6619 \\
7 & 0.135 & 0.555 & 0.7251 & 0.6534 & 0.6520 \\
8 & 0.113 & 0.538 & 0.7306 & 0.6667 & 0.6700 \\
9 & 0.091 & 0.695 & 0.7286 & 0.6345 & 0.6512 \\
10 & 0.077 & 0.736 & 0.7301 & 0.6126 & 0.6516 \\
11 & 0.062 & 0.820 & 0.7294 & 0.6000 & 0.6429 \\
12 & 0.053 & 0.890 & 0.7305 & 0.6233 & 0.6600 \\
13 & 0.047 & 0.986 & 0.7221 & 0.5701 & 0.6242 \\
14 & 0.039 & 0.953 & 0.7286 & 0.6154 & 0.6554 \\
15 & 0.036 & 1.020 & 0.7247 & 0.5945 & 0.6418 \\
16 & 0.029 & 1.161 & 0.7258 & 0.5841 & 0.6365 \\
17 & 0.028 & 1.158 & 0.7268 & 0.5874 & 0.6387 \\
18 & 0.029 & 1.178 & 0.7274 & 0.5855 & 0.6383 \\
19 & 0.025 & 1.184 & 0.7269 & 0.5808 & 0.6342 \\
20 & 0.028 & 1.182 & 0.7271 & 0.5794 & 0.6324 \\
\bottomrule
\end{tabular}
\label{tab:phase1extepochs}
\end{table*}

\subsection{Interim Findings}
\begin{itemize}[leftmargin=*]
  \item \textbf{Consistent gain from balanced sampling:} the extension improves AUROC by +0.0229 over baseline.
  \item \textbf{Best operating point appears early:} AUROC peaks around epoch 8 in the extension, after which validation loss rises sharply.
  \item \textbf{Ablation confirms fusion design:} explicit cross-attention outperforms simple concatenation by +0.0399 AUROC.
\end{itemize}

\subsection{Qualitative Error Analysis (Interim)}
We manually inspected representative false positives and false negatives from dev predictions and grouped them into recurring error modes.

\begin{table}[h]
\centering
\small
\caption{Interim failure taxonomy from manual inspection.}
\begin{tabularx}{\columnwidth}{lX}
\toprule
Category & Typical pattern \\
\midrule
Sarcasm / irony ambiguity & Surface text appears harmless but implied intent is hateful in cultural context. \\
Context beyond meme frame & External political/cultural references not visible in image-text pair. \\
Text distortion / stylization & Deliberate misspellings, split words, and symbol substitutions. \\
Entity confusion & Model latches on identity words without full relational context. \\
Image cue over-reliance & Salient visual feature dominates even when text negates hateful intent. \\
\bottomrule
\end{tabularx}
\label{tab:failuretax}
\end{table}

These observations guide planned Phase 2 work: OCR robustness, stronger context adaptation (LoRA), and abstention plus human review for uncertain cases.

\section{Progress So Far}
\subsection{Completed Work}
\begin{table*}[t]
\centering
\small
\caption{Progress summary (completed, ongoing, and planned).}
\begin{tabularx}{\textwidth}{l l X}
\toprule
Work Item & Status & Notes \\
\midrule
Dataset loading and validation & Completed & Verified paths, labels, and train/dev split consistency. \\
Baseline model (frozen CLIP + cross-attention) & Completed & End-to-end training and evaluation pipeline finalized. \\
Baseline experiment (10 epochs) & Completed & Best AUROC 0.7077; best Accuracy 67.0\%. \\
Phase 1 extension (balanced sampler, 20 epochs) & Completed & Best AUROC 0.7306 at epoch 8; +0.0229 over baseline. \\
Ablation (late fusion vs cross-attention) & Completed & Cross-attention +0.0399 AUROC over late fusion. \\
Error analysis and case categorization & Ongoing & Collecting examples for text-dominant vs image-dominant failures. \\
Phase 2 safety design (human intervention policy) & Ongoing & Defining confidence thresholds and escalation flow. \\
OCR-based end-to-end inference & Planned & Replace dataset text dependence with OCR extraction at inference. \\
LoRA/PEFT adaptation experiments & Planned & Adapt selective CLIP layers without full fine-tuning. \\
Safety evaluation protocol & Planned & Measure abstention, escalation rate, and reviewer agreement. \\
\bottomrule
\end{tabularx}
\label{tab:progress}
\end{table*}

\begin{table}[h]
\centering
\small
\caption{Internal performance comparison on dev split.}
\begin{tabular}{lcc}
\toprule
Model Variant & AUROC & Accuracy \\
\midrule
Random baseline & 0.5000 & 50.0\% \\
CLIP zero-shot (cosine) & \(\approx\)0.52 & \(\approx\)52\% \\
Phase 1 baseline & 0.7077 & 67.0\% \\
Phase 1 extension & 0.7306 & 67.0\% \\
\bottomrule
\end{tabular}
\label{tab:results}
\end{table}

\begin{table}[h]
\centering
\small
\caption{Ablation result (same training recipe, 10 epochs).}
\begin{tabular}{lcc}
\toprule
Fusion Strategy & AUROC & Macro-F1 \\
\midrule
Late Fusion & 0.6682 & 0.6127 \\
Cross-Attention Fusion & 0.7081 & 0.6320 \\
\midrule
\(\Delta\) (CrossAttn - LateFusion) & +0.0399 & +0.0193 \\
\bottomrule
\end{tabular}
\label{tab:ablation}
\end{table}

\textbf{Interpretation.}
The project has passed key mid-stage milestones: (i) stable multimodal baseline, (ii) measurable gain from class-balanced training, and (iii) architectural validation through ablation. The next bottleneck is not training stability but deployment realism and risk-sensitive moderation quality.

\subsection{Ongoing Work}
Current work focuses on three tracks:
\begin{itemize}[leftmargin=*]
  \item \textbf{Failure taxonomy:} identify recurring failure modes (sarcasm, OCR-like text distortions, rare cultural context).
  \item \textbf{Threshold policy design:} define operating thresholds for automatic accept/reject and \emph{escalate-to-human} cases.
  \item \textbf{Implementation cleanup:} improve experiment tracking and reproducibility notes for final submission.
\end{itemize}

\subsection{Planned Phase 2: Safety + Realism}
Planned Phase 2 expands the current system with practical moderation safeguards.

\textbf{Human-in-the-loop intervention.}
For uncertain or high-risk cases, the model will abstain and escalate to manual review. We will define:
\begin{itemize}[leftmargin=*,nosep]
  \item confidence bands for auto-decision vs human escalation,
  \item reviewer guideline checklist,
  \item audit logs for traceability and post-hoc correction,
  \item feedback loop to refine thresholds and retraining datasets.
\end{itemize}

\textbf{Additional deliverables.}
\begin{itemize}[leftmargin=*,nosep]
  \item OCR pipeline for end-to-end inference.
  \item LoRA-based adaptation of selected CLIP blocks.
  \item Robustness tests with noisy overlays and spelling perturbations.
\end{itemize}

\subsection{Planned Phase 2 Experiment Matrix}
To keep Phase 2 systematic, we predefine an experiment matrix instead of ad-hoc runs.

\begin{table*}[t]
\centering
\small
\caption{Planned Phase 2 experiment matrix (design before execution).}
\begin{tabularx}{\textwidth}{l l l X}
\toprule
Track & Experiment ID & Core change & Evaluation focus \\
\midrule
OCR realism & P2-OCR-1 & OCR extraction with default parser & End-to-end AUROC vs text-field oracle \\
OCR realism & P2-OCR-2 & OCR + normalization (case, punctuation, spacing) & Robustness under noisy text rendering \\
PEFT adaptation & P2-LORA-1 & LoRA on final CLIP attention blocks & AUROC/F1 gain per trainable parameter \\
PEFT adaptation & P2-LORA-2 & LoRA rank sweep (r=4,8,16) & Cost-quality tradeoff \\
Safety policy & P2-SAFE-1 & Confidence band with abstention & Escalation rate and retained accuracy \\
Safety policy & P2-SAFE-2 & High-risk keyword trigger for human review & Harm-weighted error reduction \\
Human feedback & P2-HITL-1 & Reviewer decision logging & Agreement and correction usefulness \\
Stress testing & P2-ROB-1 & Character noise and overlay perturbations & Performance degradation profile \\
\bottomrule
\end{tabularx}
\label{tab:p2matrix}
\end{table*}

\section{Ethical AI and Safety Plan}
Because this project targets harmful content moderation, we explicitly include a safety plan in technical development.

\subsection{Human Intervention Policy}
\begin{itemize}[leftmargin=*]
  \item \textbf{Low-confidence abstention:} if model confidence falls into a predefined uncertainty band, no automatic moderation action is taken.
  \item \textbf{Manual escalation:} uncertain/high-impact samples are routed to human reviewers.
  \item \textbf{Dual-review for critical cases:} content involving protected attributes can trigger secondary review.
  \item \textbf{Auditability:} model score, threshold band, and reviewer decision are logged for later audit.
\end{itemize}

\subsection{Safety Evaluation Options}
In planned Phase 2, we will add safety-specific evaluation dimensions alongside AUROC/F1:
\begin{itemize}[leftmargin=*]
  \item \textbf{Escalation rate:} fraction of samples routed to human review.
  \item \textbf{Reviewer agreement:} consistency between independent reviewers for escalated cases.
  \item \textbf{Harm-weighted error analysis:} separate accounting of false negatives in high-risk categories.
  \item \textbf{Turnaround viability:} average time from escalation to final decision.
\end{itemize}

This framing ensures the project remains not only accurate but also safer for real moderation workflows.

\subsection{Risk Register and Mitigation}
\begin{table}[h]
\centering
\small
\caption{Current project risks and mitigations.}
\begin{tabularx}{\columnwidth}{lX}
\toprule
Risk & Mitigation \\
\midrule
Overfitting after peak epoch & Early-stop checkpoints by AUROC and restrict deployment to validated windows. \\
OCR extraction noise & Add text normalization and evaluate noisy synthetic overlays before rollout. \\
Unsafe false negatives & Introduce abstention and mandatory manual review for uncertainty/high-risk bands. \\
Compute budget limits & Use PEFT (LoRA) instead of full backbone fine-tuning. \\
Annotation subjectivity & Reviewer guidelines + disagreement resolution + audit logs. \\
\bottomrule
\end{tabularx}
\label{tab:risks}
\end{table}

\section{Success Criteria and Final Evaluation Plan}
To avoid ambiguous project completion criteria, we define quantitative and process-level targets in advance.

\subsection{Quantitative Targets}
\begin{itemize}[leftmargin=*]
  \item \textbf{Primary ranking quality target:} improve or maintain AUROC relative to the current Phase 1 extension benchmark (0.7306) after introducing OCR and safety policies.
  \item \textbf{Balanced performance target:} maintain Macro-F1 close to AUROC improvements so that gains are not caused by one-class bias.
  \item \textbf{Safety-operational target:} ensure uncertain/high-risk routing does not collapse coverage (i.e., escalation rate remains practical for human moderation bandwidth).
\end{itemize}

\subsection{Operational and Reporting Targets}
\begin{itemize}[leftmargin=*]
  \item \textbf{Reproducibility target:} every table in the final report must map to a runnable command and a saved configuration.
  \item \textbf{Traceability target:} reviewer interventions for escalated cases should be logged with reasons and outcome labels.
  \item \textbf{Error-analysis target:} final write-up should include at least one validated mitigation for each major failure category in Table~\ref{tab:failuretax}.
\end{itemize}

\begin{table*}[t]
\centering
\small
\caption{Final evaluation checklist aligned with project objectives.}
\begin{tabularx}{\textwidth}{l l X}
\toprule
Objective & Acceptance signal & Evidence artifact \\
\midrule
Strong multimodal baseline & Stable re-run with near-identical trend & Training logs + checkpoint metrics \\
Class-imbalance robustness & Extension outperforms baseline on AUROC/Macro-F1 & Comparative metrics table and script output \\
Architectural validation & Cross-attention keeps advantage over late fusion & Ablation result table and run logs \\
Deployment realism & OCR-enabled path evaluated end-to-end & OCR experiment report and error examples \\
Ethical safety integration & Human-in-the-loop escalation policy operational & Escalation statistics + reviewer protocol \\
Project reproducibility & Commands and configs recreate all major results & README + config snapshots + command list \\
\bottomrule
\end{tabularx}
\label{tab:successchecklist}
\end{table*}

\section{Team Organization and Execution Workflow}
We split work by technical tracks while keeping weekly integration points to avoid drift between modeling, data processing, and safety objectives.

\begin{table}[h]
\centering
\small
\caption{Working allocation for the remaining project phase.}
\begin{tabularx}{\columnwidth}{lX}
\toprule
Track owner(s) & Responsibility \\
\midrule
Modeling track & LoRA/PEFT experiments, training stability checks, ablation updates. \\
Data/OCR track & OCR pipeline integration, normalization, and noisy text stress testing. \\
Safety track & Confidence thresholds, escalation logic, reviewer guideline sheet, audit logs. \\
Evaluation/report track & Metrics consolidation, error analysis write-up, timeline tracking. \\
\bottomrule
\end{tabularx}
\label{tab:roles}
\end{table}

\textbf{Execution protocol.}
\begin{itemize}[leftmargin=*]
  \item Weekly experiment freeze: only validated runs enter report tables.
  \item Shared artifact naming convention: run ID + config hash + timestamp.
  \item Cross-track review before merge: modeling updates and safety policy updates are reviewed together to avoid metric-only optimization.
\end{itemize}

This workflow is designed to keep the final submission balanced across model quality, safety, and reproducibility instead of optimizing a single headline metric.

\section{Tentative Timeline}
\begin{table}[h]
\centering
\small
\caption{Tentative timeline from mid-submission to final evaluation.}
\begin{tabularx}{\columnwidth}{lX}
\toprule
Time Window & Planned Deliverables \\
\midrule
Mar 9--Mar 16 & Finalize failure taxonomy; freeze Phase 2 safety policy draft. \\
Mar 17--Mar 24 & Implement OCR pipeline and confidence-based escalation logic. \\
Mar 25--Apr 2 & Run LoRA/PEFT experiments; compare with frozen-backbone extension. \\
Apr 3--Apr 10 & Safety evaluation: escalation rate, reviewer agreement, error-cost analysis. \\
Apr 11--Apr 18 & Consolidate results, finalize report/figures, prepare viva notes. \\
\bottomrule
\end{tabularx}
\label{tab:timeline}
\end{table}

\section{Repository and Reproducibility}
The public repository is:
\begin{center}
\url{https://github.com/MuditGupta2502/Meta-Hateful-Meme-Detection}
\end{center}

Main reproducibility commands used in this mid phase are:
\begin{itemize}[leftmargin=*]
  \item \texttt{python train.py --epochs 10 --batch\_size 32 --lr 1e-4}
  \item \texttt{python train.py --config configs/config\_smote.yaml}
  \item \texttt{python train\_ablation.py --config configs/config.yaml --epochs 10}
  \item \texttt{python train\_ablation.py --variants late\_fusion cross\_attention}
\end{itemize}

\section{Conclusion}
This mid-submission confirms that our project direction is viable and evidence-driven. We have already completed the core multimodal architecture and demonstrated consistent gains through a Phase 1 extension and ablation experiments. The planned Phase 2 is now focused on what the current system still lacks for realistic deployment: OCR integration, parameter-efficient adaptation, and explicit ethical safety mechanisms with human intervention on uncertain/high-risk samples.

\begin{thebibliography}{9}

\bibitem[Hu et~al., 2022]{hu2022lora}
Edward J. Hu, Yelong Shen, Phillip Wallis, Zeyuan Allen-Zhu, Yuanzhi Li,
Shean Wang, Lu Wang, and Weizhu Chen. 2022.
\newblock LoRA: Low-Rank Adaptation of Large Language Models.
\newblock In \textit{International Conference on Learning Representations (ICLR)}.

\bibitem[Kiela et~al., 2020]{kiela2020hateful}
Douwe Kiela, Hamed Firooz, Aravind Mohan, Vedanuj Goswami,
Amanpreet Singh, Pratik Ringshia, and Davide Testuggine. 2020.
\newblock The Hateful Memes Challenge: Detecting Hate Speech in Multimodal Memes.
\newblock In \textit{Advances in Neural Information Processing Systems (NeurIPS)}.

\bibitem[Lin et~al., 2017]{lin2017focal}
Tsung-Yi Lin, Priya Goyal, Ross Girshick, Kaiming He, and Piotr Dollar. 2017.
\newblock Focal Loss for Dense Object Detection.
\newblock In \textit{IEEE International Conference on Computer Vision (ICCV)}.

\bibitem[Radford et~al., 2021]{radford2021clip}
Alec Radford, Jong Wook Kim, Chris Hallacy, Aditya Ramesh, Gabriel Goh,
Sandhini Agarwal, Girish Sastry, Amanda Askell, Pamela Mishkin,
Jack Clark, Gretchen Krueger, and Ilya Sutskever. 2021.
\newblock Learning Transferable Visual Models From Natural Language Supervision.
\newblock In \textit{International Conference on Machine Learning (ICML)}.

\bibitem[Vaswani et~al., 2017]{vaswani2017attention}
Ashish Vaswani, Noam Shazeer, Niki Parmar, Jakob Uszkoreit,
Llion Jones, Aidan N. Gomez, Lukasz Kaiser, and Illia Polosukhin. 2017.
\newblock Attention Is All You Need.
\newblock In \textit{Advances in Neural Information Processing Systems (NeurIPS)}.

\bibitem[Weidinger et~al., 2021]{weidinger2021ethical}
Laura Weidinger, John Mellor, Maribeth Rauh, Conor Griffin,
Jonathan Uesato, Po-Sen Huang, Myra Cheng, and others. 2021.
\newblock Ethical and Social Risks of Harm from Language Models.
\newblock \textit{arXiv preprint arXiv:2112.04359}.

\end{thebibliography}

\end{document}
